\begin{abstract}
Architecture comparisons often use parameter counts or FLOPs as proxies for speed, but deployed systems are judged
by latency. We benchmark six small neural architectures, three single-paradigm models (sequential, parallel, modular)
and three hybrids combining depth, breadth, and modularity, all held to 385 parameters ($\pm$0.5\%) and evaluated on
three regression tasks with five seeds. A pilot study suggested a \pilotSpeedup$\times$ speedup for
modular networks. An examination traced it to a cross-framework confound and uninformative inputs, and we present it as
a case study in benchmarking pitfalls. In the corrected study, with standardized targets, every topology reaches the noise floor on simple
tasks. On a composition task, a staged modular hybrid lowers test error by \chainHBGain\% relative to a
sequential baseline ($p<0.05$), a gap that persists after 300 epochs, at
\latOneRatioHB$\times$ the single-sample latency. Across models, single-sample latency is almost perfectly predicted by tensor operations per forward call
($R^2=\opsFitRsq$) rather than by parameter count. Compilation narrows this spread (\eagerSpread$\times$ to
\compSpread$\times$) but adds a fixed per-call cost. At large
batch sizes the ranking inverts in favor of
narrow-activation designs. Small models should therefore be compared at the deployment batch size,
with operator counts reported.
\end{abstract}

\begin{IEEEkeywords}
neural network architecture, hybrid architectures, modular networks, inference latency, benchmarking,
efficient machine learning, reproducibility
\end{IEEEkeywords}
